% ===== 符号说明 =====
% 规范：自然段落形式，不用表格。每个符号一个自然段，全文统一。
\section{符号说明}

本文使用的主要数学符号及其含义按出现顺序说明如下。

% --- 示例：核心常量 ---
$N$ 表示样本总数，无量纲。该符号贯穿全文各问题，首次出现于式\cref{eq:obj}。

$\alpha$ 表示学习率，无量纲，取值范围为$(0, 1]$。在问题2的梯度下降算法中首次使用，见式\cref{eq:lr}。

% --- 示例：决策变量 ---
$x_i$ 表示第$i$个决策单元的投入量，单位为吨（t）。首次出现于问题1的目标函数，见式\cref{eq:obj}。

% --- 示例：集合与指标 ---
$\mathcal{I}$ 表示全部样本的指标集合，$\mathcal{I} = \{1, 2, \dots, N\}$。首次出现于问题1，见式\cref{eq:obj}。

% ═══════════════════════════════════════════════════════════
% 写作说明（编译前删除以下注释块）：
%
% 1. 每个符号一个自然段，格式：符号 → 含义 → 量纲/单位 → 首次出现位置（\cref引用）
% 2. 相关符号可合并在同一段落中说明（如 $x_i$ 和 $y_i$ 同属一组决策变量）
% 3. 符号按首次出现顺序排列，分批分组，读起来像连续叙述
% 4. 符号命名：优先英文字母缩写（$v$→速度、$\theta$→角度、$P$→功率）
% 5. 避免复杂符号：禁止三层上下标、过多修饰符；尽量不超过一个上下标
% 6. 全文符号统一：同一物理量在各问题中必须用同一符号
% 7. 豁免：仅在一处使用且于首次出现时已充分说明的符号可省略；
%         极常见符号（$\sum$、$\int$、$\pi$）无需说明；
%         下标约定（如 $i$ 表示指标、$t$ 表示时间）在首次出现段落中一次性交代即可
% ═══════════════════════════════════════════════════════════
